\section{Adattárolás és exportálás}

Az alkalmazásban az adattárolás szigorúan a mérés lezárásához kapcsolódik. A munkamenet közben a
rendszer a szenzoradatokat a memóriában rakja. Ezek az adatsorok csak abban az esetben kerülnek
véglegesen mentésre, ha a felhasználó szabályosan végrehajtja és befejezi a mérést. Amennyiben a
munkamenet megszakításra kerül, az alkalmazás eldobja a memóriatartalmat,

\subsection{DTO tervezési minta és Szerializáció}

A mentés megkezdése előtt az alkalmazás összeállítja a teljes munkamenetet leíró adatszerkezetet.
Az alkalmazás belső, futásidejű állapotát
reprezentáló modellek (például az aktív \texttt{SensoriaData} objektumok) nem egyeznek meg azzal
a formátummal, amely hosszú távú tárolásra optimalizált.

A probléma feloldására a rendszer a Data Transfer Object (DTO) tervezési mintát alkalmaztam.
Az aktuális állapotból egy külön exportobjektum-hierarchia (\texttt{SessionExport}, \texttt{SessionInfo},
\texttt{ExportedSensorData}) épül fel. Így mentett munkamenet önmagában is, kontextusával együtt
(metaadatok, feladatok paraméterei, normalizált gyorsulás) értelmezhető legyen a későbbi feldolgozás során.

A memóriában lévő objektumok fájlba írását a \texttt{kotlinx.serialization} keretrendszer végzi,
amely platformfüggetlen módon alakítja át a Kotlin adatosztályokat olvasásra
alkalmas  JSON (JavaScript Object Notation) szöveges formátummá \cite{kotlinserialization}.

\subsection{Biztonságos fájlkezelés}

A generált JSON fájl tárolása az Android operációs rendszer belső, izolált tárhelyén
(\texttt{Context.filesDir}) történik \cite{androidstorage}. Az itt tárolt fájlokhoz más alkalmazások
vagy külső felhasználók nem férhetnek hozzá, és az alkalmazás eltávolításakor a rendszer
automatikusan megsemmisíti őket, megelőzve az érzékeny adatok (pl. páciens azonosítók, egészségügyi mutatók)
kiszivárgását.

A fájlnevek generálása egy szigorú mintát követ,
amely tartalmazza a felhasználó nevét(Patient ID) és a mérés pontos időbélyegét
(pl. \texttt{session\_PT-001\_20260523\_...json}). Ez az elnevezés biztosítja a mérések
egyértelmű azonosíthatóságát és a mentett adatok listázását a
felhasználói felületen.

